% !TEX root = ../main.tex
\section{Ejercicios resueltos}
Los tres ejercicios se ejecutan y verifican con aserciones en \path{python/ejercicios_resueltos.py}; las salidas siguientes se generan desde ese código.

\subsection{Resuelto 1: división de prueba y Fermat}
La división llega hasta \texttt{isqrt(n)}; Fermat comprueba primero $\gcd(a,n)$ y después $a^{n-1}\bmod n$. Código y salida:
\begin{lstlisting}[language=Python,caption={División de prueba y Fermat: código ejecutable}]
def resuelto_1():
    """División de prueba y contraejemplo al recíproco de Fermat."""
    division_221 = trial_division(221)
    fermat_15 = fermat_round(15, 2)
    fermat_341 = fermat_round(341, 2)
    assert division_221 is False and fermat_15 is False and fermat_341 is True
    return {"221_es_primo": division_221, "15_pasa_fermat": fermat_15,
            "341_pasa_fermat": fermat_341}
\end{lstlisting}
\noindent\textbf{Resultado de la ejecución:}
\begin{lstlisting}[style=consolestyle]
>>> resuelto_1()
{'221_es_primo': False, '15_pasa_fermat': False, '341_pasa_fermat': True}
\end{lstlisting}

$341=11\cdot31$ pasa Fermat base $2$ pese a ser compuesto. Para $561$, la base $3$ revela un factor común; la base $2$ pasa Fermat.

\subsection{Resuelto 2: Miller--Rabin con traza verificable}
Para $221-1=220=2^2\cdot55$, la base $2$ produce:
\begin{equation}
2^{55}\bmod221=128,\qquad128^2\bmod221=30\neq220.
\end{equation}
No aparece $-1\equiv220$: $221$ es compuesto. La salida incluye parámetros, residuos, motivo y operaciones:
\begin{lstlisting}[language=Python,caption={Miller--Rabin con traza: código ejecutable}]
def resuelto_2():
    """Traza verificable de Miller--Rabin para n=221, base a=2."""
    assert decompose(221) == (2, 55)
    passed, detail = miller_rabin_round(221, 2, trace=True)
    assert not passed and detail["residuos"] == [128, 30]
    assert (detail["n"], detail["a"], detail["s"], detail["d"]) == (221, 2, 2, 55)
    return {"pasa": passed, **detail}
\end{lstlisting}
\noindent\textbf{Resultado de la ejecución:}
\begin{lstlisting}[style=consolestyle]
>>> resuelto_2()
{'pasa': False,
 'n': 221,
 'a': 2,
 's': 2,
 'd': 55,
 'residuos': [128, 30],
 'cuadrados_cadena': 1,
 'cuadrados_modulares': 6,
 'productos_modulares': 5,
 'motivo': 'sin_menos_uno'}
\end{lstlisting}

Las trazas de $341$ y $561$ se presentan junto con $221$ en la Tabla \ref{tab:trazas_propuesto1}; el registro completo está en \path{python/results/trazas_miller_rabin.csv}.

\subsection{Resuelto 3: prueba repetida y candidatos}
\begin{lstlisting}[language=Python,caption={Prueba repetida y candidatos: código ejecutable}]
def resuelto_3():
    """Candidato reproducible y cinco rondas sobre el primo conocido 1009."""
    rng = Random(2026)
    n = generate_candidate(32, rng)
    assert n.bit_length() == 32 and n % 2 == 1
    report = primality_report(1009, rounds=5, rng=Random(2026))
    assert report["resultado"] == "probable_primo"
    assert report["rondas_ejecutadas"] == 5 and report["cota_error"] == 4.0**-5
    return {"candidato": n, "bits": n.bit_length(), "impar": bool(n % 2),
            "1009_resultado": report["resultado"],
            "rondas_ejecutadas": report["rondas_ejecutadas"],
            "cota_error": report["cota_error"]}
\end{lstlisting}
\noindent\textbf{Resultado de la ejecución:}
\begin{lstlisting}[style=consolestyle]
>>> resuelto_3()
{'candidato': 2659099673,
 'bits': 32,
 'impar': True,
 '1009_resultado': 'probable_primo',
 'rondas_ejecutadas': 5,
 'cota_error': 0.0009765625}
\end{lstlisting}

Los bits extremos fijados a $1$ garantizan longitud e imparidad. Se prueba $1009$, mayor que el límite del precribado, para ejecutar las cinco rondas: se informa ``probable primo'' y la cota condicional $4^{-5}=1/1024$. La semilla fija reproduce el ejemplo, pero no es apta para claves.
